\documentclass[11pt,a4paper]{article}

% ---------------------------------------------------------
% PREAMBLE (stable, warning-free, Overleaf-safe)
% ---------------------------------------------------------
\usepackage[utf8]{inputenc}
\usepackage[T1]{fontenc}
\usepackage{lmodern}
\usepackage{amsmath, amssymb, amsfonts}
\usepackage{bm}
\usepackage{geometry}
\usepackage{graphicx}
\usepackage{hyperref}
\usepackage{cite}
\usepackage{authblk}
\usepackage{physics}
\usepackage{mathtools}

\geometry{margin=2.4cm}

\title{\textbf{Companion XXXIII: Assembly Bias in the Relaxation-Driven Cyclic Cosmology}}
\author{Michael Lehmann \\ \texttt{mi.lehmann@gmx.de}}
\date{April 2026}

\begin{document}
\maketitle

\begin{abstract}
Assembly bias describes the phenomenon that halos of identical mass cluster differently 
depending on their formation history. In the standard cosmological model, this effect 
emerges from the interplay between large-scale tidal fields, merger histories, and the 
stochasticity of hierarchical growth. In the Relaxation-Driven Cyclic Cosmology (RDCC), 
the situation is fundamentally altered. The relaxation-modified potential evolution 
introduces a temporal memory into the growth of structure, which reshapes the relation 
between halo formation time, environment, and clustering strength. This Companion 
develops a continuous, non-discrete description of assembly bias in RDCC, showing how 
the relaxation scale $k_{\mathrm{rel}}$ imprints itself on the temporal architecture of 
halo growth and on the spatial coherence of the cosmic web.
\end{abstract}
% ---------------------------------------------------------
\section{Introduction}

The concept of assembly bias challenges the simplest picture of structure formation. 
If halo mass were the only relevant variable, halos of identical mass would exhibit 
identical clustering properties. Yet numerical simulations and observations reveal a 
more intricate reality: halos that form earlier cluster more strongly, even when their 
masses are the same. This sensitivity to formation history indicates that the Universe 
retains a memory of the temporal sequence of collapse events.

In RDCC, this memory is not an emergent byproduct of hierarchical growth but a direct 
consequence of the relaxation-modified potential evolution. The suppression of 
small-scale growth and the enhancement of large-scale coherence introduce a temporal 
structure into the gravitational field itself. As a result, the formation time of a 
halo becomes intertwined with the relaxation scale, and assembly bias acquires a 
distinctive, theoretically predictable form.
% ---------------------------------------------------------
\section{Temporal Memory in RDCC Structure Formation}

In the standard cosmological model, the gravitational potential evolves primarily 
through the linear growth factor, and its time dependence is smooth and monotonic. 
RDCC modifies this picture by introducing a relaxation-driven component that depends 
on scale. Modes with $k \gg k_{\mathrm{rel}}$ experience a reduced potential decay, 
while modes with $k \ll k_{\mathrm{rel}}$ retain a more coherent temporal evolution.

This scale-dependent behavior means that the gravitational field carries a memory of 
its past configuration. When a halo begins to collapse, the potential it experiences 
is not solely determined by its local density contrast but also by the relaxation 
history of the surrounding modes. Halos forming in regions dominated by large-scale 
coherence inherit a smoother collapse trajectory, while halos forming in regions where 
small-scale modes are suppressed experience a delayed or softened collapse.

This temporal memory is the foundation of RDCC assembly bias.
% ---------------------------------------------------------
\section{Formation Time and Environmental Dependence}

In RDCC, the formation time of a halo is not an isolated parameter but a reflection of 
its environment across multiple scales. Halos embedded in regions where the potential 
evolves coherently tend to form earlier, even if their local density contrasts are 
modest. Conversely, halos in environments dominated by suppressed small-scale modes 
form later, as the relaxation mechanism slows the growth of perturbations below 
$k_{\mathrm{rel}}$.

This interplay produces a continuous spectrum of formation histories. Early-forming 
halos emerge from regions where the relaxation scale reinforces the coherence of the 
surrounding potential, while late-forming halos arise in environments where the 
relaxation mechanism inhibits the collapse of small-scale fluctuations. The result is 
a natural, scale-dependent modulation of halo formation times that directly feeds into 
assembly bias.
% ---------------------------------------------------------
\section{Clustering Signatures of RDCC Assembly Bias}

The clustering of halos reflects the coherence of the gravitational field in which 
they formed. In RDCC, early-forming halos inherit the imprint of large-scale 
relaxation, which enhances their spatial coherence. Their clustering amplitude is 
therefore elevated relative to halos of the same mass that formed later. This effect 
is not a discrete bifurcation but a smooth gradient across formation times.

Late-forming halos, shaped by environments where small-scale modes were suppressed, 
exhibit weaker clustering. Their spatial distribution reflects the diminished 
influence of coherent large-scale potentials during their collapse. The contrast 
between early and late formation is thus amplified by the relaxation scale, which 
acts as a temporal filter on the gravitational field.

The resulting assembly bias is a continuous, scale-dependent phenomenon that emerges 
naturally from the RDCC framework.
% ---------------------------------------------------------
\section{Implications for the Cosmic Web}

Assembly bias in RDCC extends beyond individual halos and influences the morphology 
of the cosmic web. Regions dominated by early-forming halos tend to develop more 
pronounced filaments and stronger connectivity, reflecting the coherence of the 
underlying potential. In contrast, regions populated by late-forming halos exhibit 
a more diffuse structure, with weaker filamentary alignment and larger voids.

This environmental dependence creates a subtle but measurable modulation of the 
cosmic web. The relaxation scale determines the degree to which large-scale modes 
shape the web's architecture, while the suppression of small-scale modes influences 
the fragmentation of filaments and the depth of voids. The interplay between these 
effects provides a rich observational landscape for testing RDCC.
% ---------------------------------------------------------
\section{Conclusion}

Assembly bias in the Relaxation-Driven Cyclic Cosmology arises from the temporal 
structure of the gravitational potential. The relaxation-modified evolution of 
perturbations introduces a memory into the collapse process, linking halo formation 
times to the coherence of their environments. This connection produces a continuous, 
scale-dependent form of assembly bias that reshapes halo clustering and the 
architecture of the cosmic web.

The results of this Companion establish assembly bias as a natural and predictive 
feature of RDCC. They also provide a conceptual bridge between the non-linear halo 
model of Companion~XXXI and future studies of galaxy formation, cluster astrophysics, 
and cosmic web morphology.

% ========================
% References
% ========================

\begin{thebibliography}{10}

\bibitem{Flagship} Lehmann, M. (2026). Relaxation-Driven Cyclic Cosmology (RDCC): A Minimal CPT-Symmetric Framework with a Single Parameter. Flagship Paper. Zenodo: \url{https://zenodo.org/records/19744958} (v43.0 or later).

\bibitem{Zenodo} The complete RDCC companion series (I--LVII) is available at the same Zenodo record above.

% Thematisch relevante Companions
\bibitem{CompXVI} Companion XVI: Boltzmann Code and Modified Hierarchy.
\bibitem{CompXXXI} Companion XXXI: RDCC Halo Model and Non-Linear Structure.
\bibitem{CompXXXII} Companion XXXII--XXXIX: Galaxy--Halo Connection, Cosmic Web, Lensing, RSD, ISW, tSZ, Rotation Curves, CGM.

% Wichtige externe Referenzen
\bibitem{Planck2020} Planck Collaboration (2020), Astron. Astrophys. 641, A6.
\bibitem{DESI2024} DESI Collaboration (2024/2025), arXiv: [aktuelle $f\sigma_8$/BAO-Paper].
\bibitem{JWST2024} Maiolino et al. (2024), Nature (JWST high-$z$ galaxies/quasars).
\bibitem{Kaiser1987} Kaiser, N. (1987), Mon. Not. R. Astron. Soc. 227, 1 (RSD).
\bibitem{Bartelmann2001} Bartelmann, M. \& Schneider, P. (2001), Phys. Rep. 340, 291 (Weak Lensing Review).
\bibitem{HuOkamoto2002} Hu, W. \& Okamoto, T. (2002), Astrophys. J. 574, 566 (CMB Lensing).
\bibitem{LewisChallinor2006} Lewis, A. \& Challinor, A. (2006), Phys. Rep. 429, 1 (CMB Lensing Review).
\bibitem{NFW1997} Navarro, J. F., Frenk, C. S. \& White, S. D. M. (1997), Astrophys. J. 490, 493 (Halo Profiles).

\end{thebibliography}

\end{document}
